\section{NHỊ THỨC NIU - TƠN}
\subsection{LÝ THUYẾT CẦN NHỚ}
\subsubsection{Nhị thức Niu-tơn}
\begin{itemize}
\item [\ding{172}] Nhắc lại các hằng đẳng thức
		\begin{boxkn}
		\begin{listEX}[2]
			\item [\ding{172}] $(a+b)^2=a^2+2ab+b^2$.
			\item [\ding{173}] $(a-b)^2=a^2-2ab+b^2$.
			\item [\ding{174}] $(a+b)^3=a^3+3a^2b+3ab^2+b^3$.
			\item [\ding{175}] $(a-b)^3=a^3-3a^2b+3ab^2-b^3$.
		\end{listEX}
	\end{boxkn}
\item [\ding{173}] Công thức nhị thức Niu - tơn
	\begin{boxkn}
		\begin{align*}
		\left(a+b\right)^n &=\mathrm{C}_n^0a^n+\mathrm{C}_n^1{a}^{n-1}b+\cdots+\mathrm{C}_n^{n-1}a{b}^{n-1}+\mathrm{C}_n^nb^n \\ 
		&=\displaystyle \sum\limits_{k=0}^n{\mathrm{C}_n^k{a}^{n-k}b^k.}
		\end{align*}
	\end{boxkn}
	\begin{itemize}
		\item [$\bullet$] Với $a=b=1$, ta có \quad $\mathrm{C}_n^0+\mathrm{C}_n^1+\cdots+\mathrm{C}_n^{n-1}+\mathrm{C}_n^n=2^n$.
		\item [$\bullet$] Với $a=1;  b=-1$, ta có \quad $\mathrm{C}_n^0-\mathrm{C}_n^1+\cdots+{\left(-1\right)}^k\mathrm{C}_n^k+\cdots+{\left(-1\right)}^n\mathrm{C}_n^n=0$.
	\end{itemize}
	\end{itemize}
\subsubsection{Chú ý}
Trong biểu thức ở vế phải của khai triển $\left(a+b\right)^n$ thì
	\begin{itemize}
		\item Số các hạng tử là $n+1$;
		\item Các hạng tử có số mũ của $a$ giảm dần từ $n$ đến $0$; số mũ của $b$ tăng dần từ $0$ đến $n$, nhưng tổng các số mũ của $a$ và $b$ trong mỗi hạng tử luôn bằng $n$ (quy ước $a^0=b^0=1$);
		\item Các hệ số của mỗi cặp hạng tử cách đều hai hạng tử đầu và cuối đều bằng nhau.
	\end{itemize}

\subsection{PHÂN LOẠI VÀ PHƯƠNG PHÁP GIẢI TOÁN}
\begin{dang}{Khai triển nhị thức Newton.} %Dạng 1
Áp dụng công thức nhị thức Newton để khai triển các biểu thức.
\end{dang}

\begin{vd} Khai triển các nhị thức sau
	\begin{tasks}(3)
		\task $(x+2)^4$
		\task $(x-2)^6$
		\task $\left(2x-1\right)^5$
		\task $\left(x+2y\right)^5$
		\task $\left(x^2-\dfrac{2}{x}\right)^5$
		\task $\left(2+\dfrac{x}{2}\right)^6$
		\task $\left(x+\dfrac{2}{x}\right)^7$
		\task ${x^2}{\left(1-2x\right)^5}$ 
		\task $\left(\dfrac{1}{2x}-2x\right)^6$
	\end{tasks}
\end{vd} \dongcham{40}

\begin{dang}{Tìm hệ số (số hạng) của $x^k$ trong khai triển $P(x)$}
	\begin{itemize}
		\item [\ding{172}] \indamm{Cách 1:} Khai triển $P(x)$, từ đó trả lời hệ số (số hạng) chứa $x^k$.
		\item [\ding{173}] \indamm{Cách 2:} Sử dụng khai triển tổng quát $(a+b)^n=\displaystyle \sum\limits_{k=0}^n{\mathrm{C}_n^k{a}^{n-k}b^k}$.
		\begin{itemize}
			\item [$\bullet$] Số hạng tổng quát trong khai triển trên là $\mathrm{C}_n^k{a}^{n-k}b^k$, với $k=0,1,2,\cdots,n$.
			\item [$\bullet$] Thu gọn phần hệ số và phần biến trong công thức vừa lập.
			\item [$\bullet$] Đồng nhất lũy thừa của biến với yêu cầu đề. Từ đây, suy ra kết quả.
		\end{itemize}
	\end{itemize}
\end{dang}
	\noindent
	\begin{tcolorbox}[colframe=cyan,colback=red!3!white,boxrule=0.5mm]
		{\color{blue}\textbf{Bài mẫu: }}Tìm hệ số của số hạng chứa $x^{12}$ trong khai triển nhị thức Niutơn của $\left(x-\dfrac{1}{x^2}\right)^{18}$
	\end{tcolorbox}
	{\centerline{\color{violet}\textbf{Lời giải:}}}
	\begin{boxkn}
	\begin{itemize}
		\item [$\bullet$] Số hạng tổng quát của khai triển trên là $$\mathrm{C}_{18}^k{x}^{18-k}\left(-\dfrac{1}{x^2} \right) ^k=\mathrm{C}_{18}^k\left(-1\right)^k\cdot x^{18-3k}, \text { với } k \in \mathbb{N} \text{ và } 0 \le k \le 18.$$
		\item [$\bullet$] Số hạng chứa $x^{12}$ tương ứng với $18-3k=12\Leftrightarrow k=2$.
		\item [$\bullet$] Vậy hệ số của số hạng chứa $x^{12}$ trong khai triển là $\left(-1\right)^2C_{18}^2=153$.
	\end{itemize}
	\end{boxkn}
\begin{vd}%[1D2Y3]
	Tìm hệ số của $x^{15}$ trong khai triển $\left(3x^2-2x\right)^{10}$.
	\loigiai{Theo công thức ở trên, ta có số hạng tổng quát trong khai triển là $$\mathrm{C}_{10}^k(3x^2)^{10-k}(-2x)^k=\mathrm{C}_{10}^k3^{10-k}(-2)^kx^{20-k}$$
		Số hạng chứa $x^{15}$ tương ứng với $20-k=15\Leftrightarrow k=5$.\\
		Vậy hệ số chứa $x^{15}$ trong khai triển là $\mathrm{C}_{10}^53^{5}(-2)^5=-6^5\cdot \mathrm{C}_{10}^5$.
	}
\end{vd} \dongcham{9}
\begin{vd}%[1D2K3]
	Tìm số hạng chứa $x^{33}$ trong các khai triển $\left(x^3+2xy\right)^{21}$.
	\loigiai{
		Số hạng tổng quát trong khai triển $\left(x^3+xy\right)^{21}$ là $$\mathrm{C}_{21}^k\left(x^3\right)^{21-k}\left(2xy\right)^k=2^k\mathrm{C}_{21}^kx^{63-2k}y^k$$
		Số hạng chứa $x^{33}$ tương ứng với số hạng có $k$ thỏa mãn $63-2k=33\Leftrightarrow k=15$.\\
		Vậy số hạng chứa $x^{23}$ trong khai triển đã cho là $2^{15}\mathrm{C}_{21}^{15}x^{33}y^{15}.$
	}
\end{vd} \dongcham{9}
\begin{vd}%01
	Tìm số hạng không chứa $x$ trong khai triển $\left(2x - \dfrac{1}{x^2}\right)^6$, với $x\ne 0$. 
	\loigiai{
		Số hạng tổng quát trong khai triển là 
		\begin{equation*}
		\mathrm{C}^k_6 (2x)^{6-k} \left(- \dfrac{1}{x^2}\right)^k = \mathrm{C}^k_6 2^{6-k} (-1)^k x^{6-3k}. 
		\end{equation*}
		
		Số hạng không chứa $x$ tương ứng với $6-3k =0 \Leftrightarrow k=2$. 
		
		Vậy số hạng cần tìm là $\mathrm{C}^2_6 2^{6-2} (-1)^2 = 240$. 
	}
\end{vd} \dongcham{11}

\begin{vd}
	Tìm hệ số $x^5$ trong khai triển $P(x)=x(1-2x)^{5} + x^2 (1+3x)^{10}$ thành đa thức.
	\loigiai{
		
		Ta có
			\begin{eqnarray*}
			x(1-2x)^{5} + x^2 (1+3x)^{10}
			&=&x \sum\limits_{k=0}^{5} \mathrm{C}_{5}^{k} (-2x)^{k}+x^{2} \sum\limits_{p=0}^{10} \mathrm{C}_{10}^{p} (3x)^{p}\\
			&=&
			\sum\limits_{k=0}^{5} \mathrm{C}_{5}^{k} \cdot x \cdot  (-2)^{k} \cdot x^{k} + \sum\limits_{p=0}^{10} \mathrm{C}_{10}^{p} \cdot x^2 \cdot  3^{p} \cdot x^{p}\\
			&=& \sum\limits_{k=0}^{5} \mathrm{C}_{5}^{k} \cdot  (-2)^{k} \cdot x^{k+1} + \sum\limits_{p=0}^{10} \mathrm{C}_{10}^{p} \cdot  3^{p} \cdot x^{p+2}.
		\end{eqnarray*}
	Suy ra hệ số của $x^5$ trong khai triển $P(x)$ ứng với $k=4; p=3$ là
		$\mathrm{C}_{5}^{4} \cdot (-2)^4 + \mathrm{C}_{10}^{3} \cdot 3^3 = 3320$.
	}
\end{vd} \dongcham{13}

\begin{vd}%[1D2B3]
	Tìm hệ số của $x^{12}$ trong khai triển $\left(x^2+1\right)^n$, biết tổng tất cả các hệ số trong khai triển đó bằng $1024$.
	\loigiai{Ta có
		$$\left(x^2+1\right)^n=\sum\limits_{k=0}^n\mathrm{C}_n^k(x^2)^k$$
		Tổng tất cả các hệ số trong khai triển là $\sum\limits_{k=0}^n\mathrm{C}_n^k=(1+1)^2=2^n$. Do đó $2^n=1024\Rightarrow n=10$.\\
		Số hạng tổng quát trong khai triển $\left(x^2+1\right)^{10}$ là $\mathrm{C}_{10}^k(x^2)^k=\mathrm{C}_{10}^kx^{2k}$.\\
		Số hạng chứa $x^{12}$ tương ứng với $2k=12\Leftrightarrow k=6$.\\
		Vậy hệ số của $x^{12}$ trong khai triển là $\mathrm{C}_{10}^6=210$.
	}
\end{vd} \dongcham{12}

\begin{vd} 
	Khai triển $(1-2x)^n=a_0+a_1 x+a_2 x^2 +\cdots +a_n x^n$. Tìm $a_5$ biết $a_0+a_1+a_2 = 71$.
	\loigiai{
		Ta có:	\begin{eqnarray*}
			\left(1-2x\right)^{n}&=&\sum\limits_{k=0}^{n} \mathrm{C}_{n}^{k} \cdot 1^{n-k} \cdot \left(-2x\right)^{k}\\
			&=&\sum\limits_{k=0}^{n} \mathrm{C}_{n}^{k}\cdot (-2)^{k} \cdot x^{k}
		\end{eqnarray*}
		Dạng tổng quát của hệ số là $a_k=\mathrm{C}_{n}^{k} (-2)^k$.
		\begin{itemize}
			\item Với $k=0 \Rightarrow a_0 = \mathrm{C}_{n}^{0} (-2)^0=1$.
			\item Với $k=1 \Rightarrow a_1 =-2 \mathrm{C}_{n}^{1}=-2n$.
			\item Với $k=2 \Rightarrow a_2 = 4 \mathrm{C}_{n}^{2} $.
		\end{itemize}
		Suy ra
			\begin{eqnarray*}
			a_0 + a_1 + a_2= 71
			&\Leftrightarrow& 1-2n+4 \mathrm{C}_{n}^{2}=71\\
			&\Leftrightarrow& 1-2n +4 \dfrac{n(n-1)}{2}=71 \\
			&\Leftrightarrow& 1-2n+2n^2-2n=71 \\
			&\Leftrightarrow& n^2 -2n - 35=0 \Leftrightarrow n=7.
		\end{eqnarray*}
		Với $n=7$ ta có hệ số của $x^5$  trong khai triển $(1-2x)^n$ là
		$$a_5=\mathrm{C}_{7}^{5} \cdot (-2)^5 = -672.$$
	}
\end{vd} \dongcham{13}

\begin{vd}
	Tìm hệ số của số hạng chứa $x^{4}$ trong khai triển $ (1+x+3x^2)^{10}$.
	\loigiai{
		Ta có: $\begin{aligned}[t]
			(1+x+3x^{2})^{10} &= [1+(x+3x^{2})]^{10} \\
			&= \sum_{k=0}^{10}  1^{10-k} \cdot C^{k}_{10} \cdot (x+3x^{2})^{k}\\
			&= \sum_{k=10}^{10}  C^{k}_{10} \sum_{p=0}^{k} \cdot C^{p}_k \cdot x^{k-p}\cdot (3x^{2})^{p}\\
			&= \sum_{k=10}^{10}  \sum_{p=0}^{k}  C^{k}_{10} \cdot C^{p}_k\cdot x^{k+p}\cdot 3^{p} 
		\end{aligned}
		$\\
		
		Số hạng chứa $x^{4}$ tương ứng với $k + p =4 $ mà $0 \le p \le k \le 10$ nên ta có các trường hợp sau
		\begin{itemize}
			\item [$\bullet$] $p=0$ thì $k=4$;
			\item [$\bullet$] $p=1$ thì $k=3$;
			\item [$\bullet$] $p=2$ thì $k=2$.
		\end{itemize}
	Vậy hệ số cần tìm là $ C^{4}_{10} \cdot C^{0}_{4} \cdot 3^{0}+ C^{3}_{10} \cdot C^{1}_{3} \cdot 3^{1} +C^{2}_{10} \cdot C^{2}_{2} \cdot 3^{2}= 1965$.
	}
\end{vd} \dongcham{13}

\subsection{BÀI TẬP TỰ LUYỆN}
\begin{baitap}
	Khai triển các biểu thức sau
	\begin{tasks}(2)
		\task $\left(x+y\right)^6$.
		\task $\left(2x-3\right)^4$.
	\end{tasks} 
	\loigiai{
		\begin{enumerate}[a)]
			\item Theo công thức nhị thức Newton ta có
			\begin{align*}
				\left(x+y\right)^6&=\mathrm{C}_6^0x^6+\mathrm{C}_6^1x^5y+\mathrm{C}_6^2x^4y^2+\mathrm{C}_6^3x^3y^3+\mathrm{C}_6^4x^2y^4+\mathrm{C}_6^5xy^5+\mathrm{C}_6^6y^6\\
				&=x^6+6x^5y+15x^4y^2+20x^3y^3+15x^2y^4+6xy^5+y^6.
			\end{align*}
			\item Theo công thức nhị thức Newton ta có
			\begin{align*}
				\left(2x-3\right)^4&=\mathrm{C}_4^0\left(2x\right)^4+\mathrm{C}_4^1\left(2x\right)^3\left(-3\right)+\mathrm{C}_4^2\left(2x\right)^2\left(-3\right)^2+\mathrm{C}_4^32x\left(-3\right)^3+\mathrm{C}_4^4\left(-3\right)^4\\
				&=16x^4-96x^3+216x^2-216x+81.
			\end{align*}
		\end{enumerate}
		
	}
\end{baitap}\cham{8}

\begin{baitap}
	Tìm hệ số của số hạng chứa $x^{6}$ trong khai triển nhị thức $(1-3x)^{11}$.
	\loigiai{
		Số hạng TQ:
		$T_{k+1} = C^{k}_{11}\cdot 1^{11-k} \cdot (-3x)^{k}= C^{k}_{11} \cdot 1^{11-k} \cdot (-3)^{k}\cdot x^{k}.$\\
		Số hạng chứa $x^{6} \Rightarrow k = 6$.\\
		Vậy hệ số cần tìm là: $C^{6}_{11} \cdot 1^{11-6}\cdot (-3)^{6}= 336798$.
	}
\end{baitap}\cham{8}

\begin{baitap}
	Tìm hệ số của số hạng chứa $x^{8}y^{9}$ trong khai triển nhị thức $ (2x-3y)^{17}$.
	\loigiai{
		Số hạng TQ:
		$\begin{aligned}[t]
			T_{k+1} &= C^{k}_{12}\cdot (2x)^{17-k}\cdot (-3y)^{k}\\
			&= C^{k}_{12}\cdot 2^{17-k}\cdot (-3)^{k}\cdot x^{17-k}\cdot y^{k}\\
		\end{aligned}$\\
		Số hạng chứa $x^{8}y^{9} \Rightarrow \heva{17-k =8 \\ k =9}\Leftrightarrow k =9$.\\
		Vậy hệ số cần tìm là: $C^{9}_{17}\cdot 2^{8}\cdot(-3)^{9}= -122494394880$.
	}
\end{baitap}\cham{8}

\begin{baitap}
	Tìm hệ số của số hạng chứa $x^{16}$ trong khai triển nhị thức $(x^{2}-2x)^{10}$.
	\loigiai{
		Số hạng TQ: 
		$\begin{aligned}[t]
			T_{k+1} &= C^{k}_{10} \cdot (x^{2})^{10-k} \cdot (-2x)^{k }\\
			&= C^{k}_{10} \cdot x^{2(10-k)} \cdot (-2)^{k} \cdot x^{k}\\
			&=C^{k}_{10} \cdot(-2)^{k}\cdot x^{20-2k} \cdot x^{k}\\
			&=C^{k}_{10} \cdot(-2)^{k}\cdot x^{20-k}.\\
		\end{aligned}$\\
		Vì có số hạng $x^{16} \Rightarrow 20 - k =16 \Leftrightarrow k =4$.\\
		Vậy hệ số cần tìm là: $C^{4}_{10}\cdot (-2)^{4} = 3360$. 
	}
\end{baitap}\cham{9}

\begin{baitap}%02
	Tìm số hạng không chứa $x$ trong khai triển $\left(x - \dfrac{2}{x^2}\right)^{21}$, với $x \ne 0$. 
	\loigiai{
		Số hạng tổng quát trong khai triển là 
		\begin{equation*}
			\mathrm{C}^k_{21} (x)^{21-k} \left(- \dfrac{2}{x^2}\right)^k = \mathrm{C}^k_{21} (-2)^{k} x^{21-3k}. 
		\end{equation*}
		
		Ta phải tìm $k$ sao cho $21-3k =0 \Leftrightarrow k=7$. 
		
		Vậy số hạng cần tìm là $\mathrm{C}^7_{21} (-2)^{7}$. 
	}
\end{baitap}\cham{9}
